% GENERATED FILE. Edit canonical lessons, then run npm run generate:books.
% canonical-source-hash: fnv1a64:14793eee19586edb
% canonical-lessons: ES-C41-creer, ES-C41-asi-que, ES-C41-deber, ES-C41-explicar

\chapter{Saying Why}
\label{ch:spanish-giving-reasons}
\hlchaptermodality{115}

\begin{chapteropening}
\textbf{By the end of this chapter:} \emph{I can give a reason for what I think and what I plan — stating an opinion with creo que, running forward to its consequence with así que, and naming an obligation with debo.}

Answer ¿Por qué no corres? in four moves — Creo que llueve. No corro porque llueve, así que leo. Debo estudiar. Opinion, cause, consequence, obligation, each with its own word, and every word in it taught by this book.
\end{chapteropening}

\section[creer]{creer — to put your heart somewhere}

\label{lesson:ES-C41-creer}



\begin{quote}
You have \textbf{pensar} — to think, and underneath it \emph{pēnsāre}, \textquotedblleft{}to weigh\textquotedblright{}. Spanish has a second verb for the inside of your head, and it does a different job. Where \emph{pensar} weighs, this one commits.
\end{quote}

\begin{sounds}
\begin{itemize}
\raggedright
  \item \texttt{hiatus-ee} — the two \emph{e}'s do not merge: \textbf{creer} = \emph{kreh-EHR}, two beats.
  \item \texttt{r-tap} — a single flick, not the rolled \emph{rr}.
\end{itemize}
\end{sounds}

\begin{cousinweb}
\textbf{creer} = \textquotedblleft{}to believe\textquotedblright{}, and in everyday use \textquotedblleft{}to think\textquotedblright{} in the sense of \emph{hold an opinion}.

It comes from Latin \emph{\textbf{crēdere}}, \textquotedblleft{}to entrust, to believe\textquotedblright{} — and \emph{crēdere} is a compound so old that Latin had already forgotten it was one. It is PIE \emph{\textbf{*\'{k}red-d\textsuperscript{h}eh\textsubscript{1}-}}: "\textbf{to put} \emph{\textbf{*\'{k}red-}}".

What is that first half? Traditionally, the \textbf{heart} word — the root of Latin \emph{\textbf{cor}}, \emph{\textbf{cordis}} and of English \textbf{heart}. Be honest that the step is \textbf{disputed}: that root normally appears as \emph{*\'{k}\'{\={e}}r}, never \emph{*\'{k}red-}. The compound is not in doubt; only what it was built on.

Taking the traditional reading:

\noindent
\begin{tabularx}{\linewidth}{@{}>{\raggedright\arraybackslash}X>{\raggedright\arraybackslash}X@{}}
\toprule
\textbf{from \emph{crēdere}} & \textbf{from \emph{cor, cordis}} \\
\midrule
\textbf{credit}, \textbf{creed}, \textbf{credo} & \textbf{cordial}, \textbf{courage} \\
\textbf{credential}, \textbf{incredible} & \textbf{record}, \textbf{accord}, \textbf{discord} \\
\bottomrule
\end{tabularx}

To \textbf{record} is to bring something back to the \emph{cor} — for a Roman the seat of \textbf{memory}, not of feeling. To have \textbf{courage} is to have heart.
\end{cousinweb}

\begin{grammarlens}[title={Grammar Lens: creer que}]
\emph{Creer} almost always drags a \textbf{que} behind it, and then a whole clause:

\noindent
\begin{tabularx}{\linewidth}{@{}>{\raggedright\arraybackslash}X>{\raggedright\arraybackslash}X@{}}
\toprule
\textbf{} & \textbf{} \\
\midrule
\textbf{Creo que} llueve. & I think it's raining. \\
\textbf{Creo que} hace frío. & I think it's cold. \\
\textbf{No creo que} llu\textbf{e}va. & I don't think it's raining. \\
\bottomrule
\end{tabularx}

Look hard at that third row. \emph{Creo que} \textbf{llueve} — but \emph{no creo que} \textbf{llueva}. Negating \emph{creer} flips the clause into the \textbf{subjunctive} from chapter eighteen: stop vouching for something and Spanish stops using the tense that states facts.

\emph{\textbf{Pienso}} is the weighing; \emph{\textbf{creo}} is where it landed. Treat that as a way to feel the difference, not a rule — both take \emph{que} and a clause, and a speaker may use either. But \emph{creo que} is the ordinary opening for an opinion.
\end{grammarlens}

\subsection*{Your turn}
Say these aloud:

\begin{itemize}
\raggedright
  \item \textquotedblleft{}creer\textquotedblright{} — \emph{kreh-EHR}, two beats, not one
  \item \textquotedblleft{}Creo que llueve.\textquotedblright{}
  \item the pair — \textquotedblleft{}pienso\textquotedblright{} (I weigh) then \textquotedblleft{}creo\textquotedblright{} (I hold)
  \item \textquotedblleft{}creer\textquotedblright{} then English \textquotedblleft{}credit, creed, cordial, record\textquotedblright{}
\end{itemize}

\begin{tcolorbox}[breakable,colback=teal!4,colframe=teal!35!black,title={Before you move on}]
What shape is \emph{crēdere}, under the surface? (\textbf{A compound} — \emph{*\'{k}red-} + \emph{d\textsuperscript{h}eh\textsubscript{1}-}, "to put \emph{*\'{k}red-}".) What is the traditional reading of that first half, and why hedge it? (\textbf{Heart} — but the heart root is \emph{*\'{k}\'{\={e}}r}, never \emph{*\'{k}red-}, so the step is disputed.) \emph{Pensar} weighs; what does \emph{creer} do? (It \textbf{lands} — it holds the opinion the weighing produced.)
\end{tcolorbox}
\section[así que]{así que — \textquotedblleft{}so\textquotedblright{}, and a word you have been saying since yes}

\label{lesson:ES-C41-asi-que}



\begin{quote}
\emph{Porque} runs from a claim back to its reason. This runs the other way — from the reason forward to what follows. Same two facts, opposite arrow, and you need both to explain anything.
\end{quote}

\begin{sounds}
\begin{itemize}
\raggedright
  \item \texttt{s-clear} — a clean hiss in \textbf{así}, never a \emph{z} buzz.
  \item \texttt{k-hard} — \emph{qu} is a plain \emph{k}; the \emph{u} is silent.
\end{itemize}
\end{sounds}

\begin{cousinweb}
\textbf{así} = \textquotedblleft{}so, thus, in that way\textquotedblright{}. \textbf{así que} = \textquotedblleft{}so, and therefore\textquotedblright{}.

\emph{Así} is Latin \emph{\textbf{sīc}}, \textquotedblleft{}thus\textquotedblright{} — with an \emph{a-} that Spanish added later, the same one it put on \emph{apenas} and \emph{afuera}. And you have met \emph{sīc} before: it is the whole of Spanish \textbf{sí}, \textquotedblleft{}yes\textquotedblright{}. Latin \emph{sīc} meant \textquotedblleft{}thus, in this way\textquotedblright{}, and saying \emph{thus} in answer to a question is how a language gets a word for yes.

So \emph{sí} and \emph{así} are the same Latin adverb, one bare and one with the \emph{a-} attached. English borrowed it raw and never translated it: \emph{\textbf{sic}}, the mark an editor puts after a quoted mistake — \textquotedblleft{}thus it stood\textquotedblright{}.
\end{cousinweb}

\begin{grammarlens}[title={Grammar Lens: the two arrows}]
The pair is worth drilling as a pair, because they carry the same two facts in opposite orders:

\noindent
\begin{tabularx}{\linewidth}{@{}>{\raggedright\arraybackslash}X>{\raggedright\arraybackslash}X@{}}
\toprule
\textbf{} & \textbf{} \\
\midrule
No corro \textbf{porque} llueve. & I'm not running \textbf{because} it's raining. \\
Llueve, \textbf{así que} no corro. & It's raining, \textbf{so} I'm not running. \\
\bottomrule
\end{tabularx}

\emph{Porque} points \textbf{backwards} to the cause. \emph{Así que} points \textbf{forwards} to the consequence. Which you reach for depends only on which fact you want to say first — and that is a choice about emphasis, not about grammar.

\noindent
\begin{tabularx}{\linewidth}{@{}>{\raggedright\arraybackslash}X>{\raggedright\arraybackslash}X@{}}
\toprule
\textbf{} & \textbf{} \\
\midrule
Hace frío, \textbf{así que} bebo café. & It's cold, so I drink coffee. \\
\textbf{Creo que} llueve, \textbf{así que} leo. & I think it's raining, so I read. \\
\bottomrule
\end{tabularx}
\end{grammarlens}

\subsection*{Your turn}
Say these aloud:

\begin{itemize}
\raggedright
  \item \textquotedblleft{}sí\textquotedblright{} then \textquotedblleft{}así\textquotedblright{} — the same Latin word, one with an \emph{a-} on the front
  \item \textquotedblleft{}Llueve, así que no corro.\textquotedblright{}
  \item the same two facts the other way — \textquotedblleft{}No corro porque llueve.\textquotedblright{}
  \item \textquotedblleft{}Creo que hace frío, así que bebo café.\textquotedblright{}
\end{itemize}

\begin{tcolorbox}[breakable,colback=teal!4,colframe=teal!35!black,title={Before you move on}]
Which Latin adverb is hiding in both \emph{sí} and \emph{así}? (\emph{\textbf{Sīc}}, \textquotedblleft{}thus\textquotedblright{}.) Which way does \emph{porque} point, and which way \emph{así que}? (\emph{Porque} \textbf{backwards} to the cause; \emph{así que} \textbf{forwards} to the consequence.) What does an editor's \emph{sic} mean? (\textbf{Thus it stood} — the same word, borrowed whole.)
\end{tcolorbox}
\section[deber]{deber — to owe, and therefore to ought}

\label{lesson:ES-C41-deber}



\begin{quote}
You can now say what you think and what follows from it. This verb adds the third move in any explanation: what somebody \textbf{ought} to do about it. It gets there from an unexpected direction — money.
\end{quote}

\begin{sounds}
\begin{itemize}
\raggedright
  \item \texttt{b-soft} — between vowels the \emph{b} softens: \textbf{deber} = \emph{deh-BEHR}.
  \item \texttt{r-tap} — one flick at the end, not a roll.
\end{itemize}
\end{sounds}

\begin{cousinweb}
\textbf{deber} does two jobs: \textquotedblleft{}\textbf{to owe}\textquotedblright{}, and \textquotedblleft{}\textbf{ought to, should}\textquotedblright{}.

Latin \emph{\textbf{dēbēre}} explains both, because it is a compound: \emph{\textbf{dē-}} (\textquotedblleft{}from\textquotedblright{}) plus \emph{\textbf{habēre}} (\textquotedblleft{}to have\textquotedblright{}). To \emph{dēbēre} something is \textbf{to have it from someone} — which is exactly what owing is. You are holding what is theirs.

English helped itself to this verb four times over — three of them from the participle \emph{dēbitum} rather than the verb — and the spellings record each route:

\noindent
\begin{tabularx}{\linewidth}{@{}>{\raggedright\arraybackslash}X>{\raggedright\arraybackslash}X@{}}
\toprule
\textbf{English} & \textbf{how it came} \\
\midrule
\textbf{debt} & through French, then re-spelled to show the Latin \emph{b} \\
\textbf{debit} & straight from Latin bookkeeping \\
\textbf{due} & through French \emph{deü}, the \emph{b} worn away \\
\textbf{duty} & built on \emph{due} \\
\bottomrule
\end{tabularx}

One more came the long way round, and not by borrowing. French turned \emph{dēbēre} into \emph{devoir}, \textquotedblleft{}duty\textquotedblright{}. English \textbf{translated} the phrase \emph{mettre en devoir} — to put oneself under an obligation — half-way, as \emph{put in dever}, and the last two words later fused into \textbf{endeavour}.
\end{cousinweb}

\begin{grammarlens}[title={Grammar Lens: debo + infinitive}]
\emph{Deber} is regular, and for \textquotedblleft{}ought to\textquotedblright{} it takes a bare infinitive straight after it — no \emph{que}, no preposition:

\noindent
\begin{tabularx}{\linewidth}{@{}>{\raggedright\arraybackslash}X>{\raggedright\arraybackslash}X@{}}
\toprule
\textbf{person} & \textbf{form} \\
\midrule
yo & \textbf{debo} \\
tú & \textbf{debes} \\
él/ella/usted & \textbf{debe} \\
nosotros & \textbf{debemos} \\
ellos/ustedes & \textbf{deben} \\
\bottomrule
\end{tabularx}

\textbf{Debo estudiar.} — I ought to study. \textbf{Llueve, así que debo leer.} — It's raining, so I ought to read.

The obligation is softer than a command and firmer than a wish, which is the register an explanation usually wants.
\end{grammarlens}

\subsection*{Your turn}
Say these aloud:

\begin{itemize}
\raggedright
  \item \textquotedblleft{}debo, debes, debe, debemos, deben\textquotedblright{}
  \item \textquotedblleft{}Debo trabajar.\textquotedblright{}
  \item \textquotedblleft{}Hace frío, así que debo beber café.\textquotedblright{}
  \item \textquotedblleft{}Creo que llueve, así que debo leer.\textquotedblright{}
  \item \textquotedblleft{}deber\textquotedblright{} then English \textquotedblleft{}debt, debit, due, duty\textquotedblright{}
\end{itemize}

\begin{tcolorbox}[breakable,colback=teal!4,colframe=teal!35!black,title={Before you move on}]
What two Latin pieces build \emph{dēbēre}, and what do they say together? (\emph{\textbf{Dē-}} + \emph{\textbf{habēre}} — \textquotedblleft{}to have \textbf{from} someone\textquotedblright{}, which is owing.) Name two of the four English words it became. (Any of \textbf{debt}, \textbf{debit}, \textbf{due}, \textbf{duty}.) What follows \emph{debo} — a \emph{que}, or a bare infinitive? (\textbf{A bare infinitive}: \emph{debo estudiar}.)
\end{tcolorbox}
\section[explicar]{explicar — to unfold}

\label{lesson:ES-C41-explicar}



\begin{quote}
Three pieces are in your hands: \emph{creo que} to state an opinion, \emph{así que} to run forward to what follows, \emph{debo} to say what ought to happen. This lesson puts all three in a row, under the verb that names the act.
\end{quote}

\begin{sounds}
\begin{itemize}
\raggedright
  \item \texttt{k-hard} — \emph{c} before \emph{a} is a hard \emph{k}: \textbf{explicar} = \emph{eks-plee-KAR}.
  \item \texttt{l-clear} — a clean \emph{l} in the \emph{pl}, no vowel slipped in between.
\end{itemize}
\end{sounds}

\begin{cousinweb}
\textbf{explicar} = Latin \emph{\textbf{explicāre}}: \emph{\textbf{ex-}} (\textquotedblleft{}out\textquotedblright{}) + \emph{\textbf{plicāre}} (\textquotedblleft{}to fold\textquotedblright{}). \textbf{To explain is to unfold} — something was folded shut, and you open it.

\emph{Plicāre} is one of the great suppliers of English. Some came from Latin directly, some through French, and the sound changed on the way:

\noindent
\begin{tabularx}{\linewidth}{@{}>{\raggedright\arraybackslash}X>{\raggedright\arraybackslash}X@{}}
\toprule
\textbf{straight from Latin} & \textbf{through French} \\
\midrule
\textbf{explicate} & \textbf{explicit}, \textbf{application} \\
\textbf{complicated} (\emph{com-} + fold) & \textbf{employ}, \textbf{deploy} \\
\textbf{duplicate} & \textbf{display} — a doublet of \emph{deploy} \\
\bottomrule
\end{tabularx}

And \textbf{simple} is \emph{sim-plex} — \textquotedblleft{}\textbf{one}-fold\textquotedblright{}. Something simple has been folded only once. Something \textbf{complicated} has been folded \emph{together}, over and over, which is why it needs unfolding. The metaphor is consistent all the way down, and Spanish keeps it in the plainest verb of the set.
\end{cousinweb}

\begin{grammarlens}[title={Grammar Lens: the shape of an explanation}]
\emph{Explicar} is regular — it takes the plain \emph{-ar} endings of \emph{hablar}, with no stem change at all: \textbf{explico}, \textbf{explicas}, \textbf{explica}, \textbf{explicamos}, \textbf{explican}. (\emph{Contar}, which you met earlier, breaks its \emph{o} to \emph{ue}; this one does not, and there is nothing to remember.)

Here is a whole explanation, with every piece something this book has taught:

\begin{quote}
— ¿Por qué no corres? — \textbf{Creo que} llueve. No corro \textbf{porque} llueve, \textbf{así que} leo. \textbf{Debo} estudiar.
\end{quote}

>

\begin{quote}
\emph{— Why aren't you running?} \emph{— I think it's raining. I'm not running because it's raining, so I read. I ought to study.}
\end{quote}

Four moves, four words: an \textbf{opinion} (\emph{creo que}), a \textbf{cause} (\emph{porque}), a \textbf{consequence} (\emph{así que}), an \textbf{obligation} (\emph{debo}). That is what an explanation is made of, and you can now build one in either direction — cause first or claim first — because \emph{porque} and \emph{así que} let you choose.
\end{grammarlens}

\subsection*{Your turn}
Say these aloud:

\begin{itemize}
\raggedright
  \item \textquotedblleft{}explico, explicas, explica, explicamos, explican\textquotedblright{}
  \item \textquotedblleft{}Explico bien.\textquotedblright{} — I explain well.
  \item the four moves — \textquotedblleft{}Creo que llueve, porque hace frío, así que debo leer.\textquotedblright{}
  \item \textquotedblleft{}explicar\textquotedblright{} then English \textquotedblleft{}explicit, complicated, simple, display\textquotedblright{}
\end{itemize}

\begin{tcolorbox}[breakable,colback=teal!4,colframe=teal!35!black,title={Before you move on}]
What does \emph{explicāre} literally say? (\textbf{To fold out} — \emph{ex-} + \emph{plicāre}.) What is \emph{simple}, folded? (\textbf{Once} — \emph{sim-plex}; \emph{complicated} is folded together, again and again.) Name the four moves of an explanation and the word for each. (Opinion \emph{\textbf{creo que}}, cause \emph{\textbf{porque}}, consequence \emph{\textbf{así que}}, obligation \emph{\textbf{debo}}.)
\end{tcolorbox}
